% ============================================================
\chapter{SQL --- Jointures et Sous-requ\^etes}
\label{ch:jointures}
% ============================================================

La puissance du mod\`ele relationnel r\'eside dans sa capacit\'e \`a relier des tables entre elles. Un \'etudiant est inscrit dans un d\'epartement, qui appartient \`a une universit\'e~; chaque inscription r\'ef\'erence un cours, qui est enseign\'e par un professeur. Pour extraire l'information compl\`ete --- par exemple, \og quels \'etudiants suivent les cours du Professeur Dupont~? \fg --- il faut \emph{joindre} plusieurs tables. La jointure est l'op\'eration reine de SQL, et sa ma\^itrise distingue le d\'ebutant de l'expert. Les sous-requ\^etes, quant \`a elles, permettent d'imbriquer des questions les unes dans les autres, offrant une expressivit\'e consid\'erable.

\section{Introduction}

Les \emph{jointures} permettent de combiner des donn\'ees provenant de plusieurs
tables en exploitant les relations entre elles (cl\'es \'etrang\`eres).
Les \emph{sous-requ\^etes} (ou requ\^etes imbriqu\'ees) permettent d'utiliser
le r\'esultat d'une requ\^ete comme entr\'ee d'une autre.

\section{Jointure interne (INNER JOIN)}

\begin{definition}[Jointure interne]
La \emph{jointure interne} combine les lignes de deux tables o\`u la condition
de jointure est satisfaite. Les lignes sans correspondance sont \'elimin\'ees.
\end{definition}

\begin{codeblock}[title=\textbf{Syntaxes \'equivalentes de la jointure interne}]
\begin{minted}{sql}
-- Syntaxe explicite (recommandee)
SELECT e.nom, c.intitule, i.note
FROM Etudiants e
INNER JOIN Inscriptions i ON e.id_etudiant = i.id_etudiant
INNER JOIN Cours c ON i.id_cours = c.id_cours;

-- Syntaxe implicite (ancienne, a eviter)
SELECT e.nom, c.intitule, i.note
FROM Etudiants e, Inscriptions i, Cours c
WHERE e.id_etudiant = i.id_etudiant
  AND i.id_cours = c.id_cours;
\end{minted}
\end{codeblock}

\begin{codeoutput}
\begin{verbatim}
nom     | intitule         | note
--------+------------------+------
Bernard | Bases de donnees | 15.5
Bernard | Algorithmes      | 13.0
Petit   | Bases de donnees | 12.0
Petit   | Reseaux          | 16.0
Moreau  | Algorithmes      | 17.5
Moreau  | Statistiques     | 14.0
Leroy   | Bases de donnees | 9.0
Roux    | Statistiques     | 18.0
\end{verbatim}
\end{codeoutput}

\begin{bonnepratique}[title=\textbf{Toujours utiliser la syntaxe explicite JOIN}]
La syntaxe implicite (virgule dans le \texttt{FROM}) est source d'erreurs~:
un oubli de condition dans le \texttt{WHERE} produit un produit cart\'esien
accidentel. La syntaxe \texttt{JOIN ... ON} est plus claire et s\'eparere
la logique de jointure du filtrage.
\end{bonnepratique}

\section{Jointure naturelle (NATURAL JOIN)}

\begin{codeblock}[title=\textbf{Jointure naturelle}]
\begin{minted}{sql}
-- Joint sur les colonnes de meme nom
SELECT * FROM Etudiants NATURAL JOIN Inscriptions;
\end{minted}
\end{codeblock}

\begin{attention}[title=\textbf{Danger de NATURAL JOIN}]
La jointure naturelle d\'etermine automatiquement les colonnes communes.
Si deux tables partagent un nom de colonne par co\"incidence (ex.\ \texttt{nom}
dans \texttt{Etudiants} et \texttt{Professeurs}), le r\'esultat sera incorrect.
Pr\'ef\'erez toujours la jointure explicite avec \texttt{ON}.
\end{attention}

\section{Jointures externes (OUTER JOIN)}

\begin{definition}[Jointure externe]
Une \emph{jointure externe} conserve toutes les lignes d'une table (ou des deux),
m\^eme celles sans correspondance dans l'autre table. Les colonnes manquantes
sont remplies par \texttt{NULL}.
\end{definition}

\subsection{LEFT JOIN}

\begin{codeblock}[title=\textbf{LEFT JOIN --- tous les \'etudiants, m\^eme sans inscription}]
\begin{minted}{sql}
SELECT e.nom, e.prenom, c.intitule, i.note
FROM Etudiants e
LEFT JOIN Inscriptions i ON e.id_etudiant = i.id_etudiant
LEFT JOIN Cours c ON i.id_cours = c.id_cours
ORDER BY e.nom;
\end{minted}
\end{codeblock}

Si un \'etudiant n'est inscrit \`a aucun cours, il appara\^it quand m\^eme avec
des \texttt{NULL} pour \texttt{intitule} et \texttt{note}.

\subsection{RIGHT JOIN}

\begin{codeblock}[title=\textbf{RIGHT JOIN --- tous les cours, m\^eme sans inscription}]
\begin{minted}{sql}
SELECT c.intitule, e.nom, i.note
FROM Inscriptions i
RIGHT JOIN Cours c ON i.id_cours = c.id_cours
LEFT JOIN Etudiants e ON i.id_etudiant = e.id_etudiant;
\end{minted}
\end{codeblock}

\begin{remarque}
\texttt{RIGHT JOIN} est moins utilis\'e car on peut toujours le remplacer
par un \texttt{LEFT JOIN} en inversant l'ordre des tables.
SQLite ne supporte pas \texttt{RIGHT JOIN} (avant la version 3.39).
\end{remarque}

\subsection{FULL OUTER JOIN}

\begin{codeblock}[title=\textbf{FULL OUTER JOIN}]
\begin{minted}{sql}
-- Tous les etudiants et tous les cours (PostgreSQL)
SELECT e.nom, c.intitule
FROM Etudiants e
FULL OUTER JOIN Inscriptions i ON e.id_etudiant = i.id_etudiant
FULL OUTER JOIN Cours c ON i.id_cours = c.id_cours;
\end{minted}
\end{codeblock}

\begin{formules}[title=\textbf{R\'esum\'e des types de jointures}]
\begin{center}
\begin{tabular}{lp{8cm}}
\toprule
\textbf{Type} & \textbf{Description} \\
\midrule
\texttt{INNER JOIN} & Lignes avec correspondance dans les deux tables \\
\texttt{LEFT JOIN} & Toutes les lignes de la table gauche + correspondances \\
\texttt{RIGHT JOIN} & Toutes les lignes de la table droite + correspondances \\
\texttt{FULL OUTER JOIN} & Toutes les lignes des deux tables \\
\texttt{CROSS JOIN} & Produit cart\'esien (toutes les combinaisons) \\
\bottomrule
\end{tabular}
\end{center}
\end{formules}

\section{Jointure crois\'ee (CROSS JOIN)}

\begin{codeblock}[title=\textbf{Produit cart\'esien}]
\begin{minted}{sql}
-- Toutes les combinaisons etudiant-cours
SELECT e.nom, c.intitule
FROM Etudiants e
CROSS JOIN Cours c;
-- 5 etudiants x 4 cours = 20 lignes
\end{minted}
\end{codeblock}

\section{Auto-jointure (Self Join)}

\begin{definition}[Auto-jointure]
Une \emph{auto-jointure} est une jointure d'une table avec elle-m\^eme.
Elle n\'ecessite l'utilisation d'alias pour distinguer les deux
\og copies \fg de la table.
\end{definition}

\begin{codeblock}[title=\textbf{Trouver les paires d'\'etudiants de la m\^eme fili\`ere}]
\begin{minted}{sql}
SELECT e1.nom AS etudiant1, e2.nom AS etudiant2, e1.filiere
FROM Etudiants e1
JOIN Etudiants e2 ON e1.filiere = e2.filiere
                  AND e1.id_etudiant < e2.id_etudiant;
\end{minted}
\end{codeblock}

\begin{codeoutput}
\begin{verbatim}
etudiant1 | etudiant2 | filiere
----------+-----------+--------
Bernard   | Petit     | Info
Bernard   | Leroy     | Info
Petit     | Leroy     | Info
Moreau    | Roux      | Maths
\end{verbatim}
\end{codeoutput}

\section{Sous-requ\^etes}

\begin{definition}[Sous-requ\^ete]
Une \emph{sous-requ\^ete} est une requ\^ete \texttt{SELECT} imbriqu\'ee
dans une autre requ\^ete. Elle peut appara\^itre dans les clauses
\texttt{WHERE}, \texttt{FROM}, ou \texttt{SELECT}.
\end{definition}

\subsection{Sous-requ\^ete scalaire}

Une sous-requ\^ete qui retourne exactement une valeur~:

\begin{codeblock}[title=\textbf{Sous-requ\^ete scalaire}]
\begin{minted}{sql}
-- Etudiants ayant une note superieure a la moyenne
SELECT e.nom, i.note
FROM Etudiants e
JOIN Inscriptions i ON e.id_etudiant = i.id_etudiant
WHERE i.note > (SELECT AVG(note) FROM Inscriptions);
\end{minted}
\end{codeblock}

\begin{codeoutput}
\begin{verbatim}
-- Moyenne = 14.375
nom     | note
--------+------
Bernard | 15.5
Petit   | 16.0
Moreau  | 17.5
Roux    | 18.0
\end{verbatim}
\end{codeoutput}

\subsection{Sous-requ\^ete avec IN}

\begin{codeblock}[title=\textbf{Sous-requ\^ete avec IN}]
\begin{minted}{sql}
-- Etudiants inscrits a un cours du professeur Dupont
SELECT nom, prenom
FROM Etudiants
WHERE id_etudiant IN (
    SELECT id_etudiant
    FROM Inscriptions
    WHERE id_cours IN (
        SELECT id_cours FROM Cours WHERE id_prof = 1
    )
);
\end{minted}
\end{codeblock}

\subsection{Sous-requ\^ete corr\'el\'ee}

\begin{definition}[Sous-requ\^ete corr\'el\'ee]
Une sous-requ\^ete est \emph{corr\'el\'ee} lorsqu'elle r\'ef\'erence une colonne
de la requ\^ete externe. Elle est r\'e\'evalu\'ee pour chaque ligne de la
requ\^ete externe.
\end{definition}

\begin{codeblock}[title=\textbf{Sous-requ\^ete corr\'el\'ee avec EXISTS}]
\begin{minted}{sql}
-- Etudiants inscrits a au moins un cours
SELECT e.nom, e.prenom
FROM Etudiants e
WHERE EXISTS (
    SELECT 1
    FROM Inscriptions i
    WHERE i.id_etudiant = e.id_etudiant
);

-- Etudiants NON inscrits a aucun cours
SELECT e.nom, e.prenom
FROM Etudiants e
WHERE NOT EXISTS (
    SELECT 1
    FROM Inscriptions i
    WHERE i.id_etudiant = e.id_etudiant
);
\end{minted}
\end{codeblock}

\begin{attention}[title=\textbf{Performance des sous-requ\^etes corr\'el\'ees}]
Les sous-requ\^etes corr\'el\'ees sont r\'e\'evalu\'ees pour chaque ligne.
Sur de grandes tables, elles peuvent \^etre lentes. Les optimiseurs
modernes les transforment souvent en jointures, mais il est bon de
conna\^itre l'alternative avec \texttt{JOIN}~:
\begin{minted}{sql}
-- Equivalent avec LEFT JOIN + IS NULL
SELECT e.nom FROM Etudiants e
LEFT JOIN Inscriptions i ON e.id_etudiant = i.id_etudiant
WHERE i.id_etudiant IS NULL;
\end{minted}
\end{attention}

\subsection{Sous-requ\^ete dans FROM (table d\'eriv\'ee)}

\begin{codeblock}[title=\textbf{Sous-requ\^ete dans FROM}]
\begin{minted}{sql}
-- Moyenne par etudiant, puis filtrage
SELECT sub.nom, sub.moyenne
FROM (
    SELECT e.nom, AVG(i.note) AS moyenne
    FROM Etudiants e
    JOIN Inscriptions i ON e.id_etudiant = i.id_etudiant
    GROUP BY e.id_etudiant, e.nom
) AS sub
WHERE sub.moyenne > 14;
\end{minted}
\end{codeblock}

\begin{codeoutput}
\begin{verbatim}
nom     | moyenne
--------+--------
Bernard | 14.25
Petit   | 14.0
Moreau  | 15.75
Roux    | 18.0
\end{verbatim}
\end{codeoutput}

\subsection{Sous-requ\^ete dans SELECT}

\begin{codeblock}[title=\textbf{Sous-requ\^ete dans SELECT}]
\begin{minted}{sql}
SELECT e.nom,
       (SELECT COUNT(*) FROM Inscriptions i
        WHERE i.id_etudiant = e.id_etudiant) AS nb_cours
FROM Etudiants e;
\end{minted}
\end{codeblock}

\begin{codeoutput}
\begin{verbatim}
nom     | nb_cours
--------+---------
Bernard | 2
Petit   | 2
Moreau  | 2
Leroy   | 1
Roux    | 1
\end{verbatim}
\end{codeoutput}

\section{Op\'erateurs ALL, ANY/SOME}

\begin{codeblock}[title=\textbf{Comparaison avec ALL et ANY}]
\begin{minted}{sql}
-- Note superieure a TOUTES les notes du cours 101
SELECT e.nom, i.note
FROM Etudiants e
JOIN Inscriptions i ON e.id_etudiant = i.id_etudiant
WHERE i.note > ALL (
    SELECT note FROM Inscriptions WHERE id_cours = 101
);

-- Note superieure a AU MOINS UNE note du cours 101
SELECT e.nom, i.note
FROM Etudiants e
JOIN Inscriptions i ON e.id_etudiant = i.id_etudiant
WHERE i.note > ANY (
    SELECT note FROM Inscriptions WHERE id_cours = 101
);
\end{minted}
\end{codeblock}

\section{Expressions de table communes (CTE)}

\begin{definition}[CTE (Common Table Expression)]
Une \emph{CTE} est une requ\^ete nomm\'ee temporaire, d\'efinie avec le mot-cl\'e
\texttt{WITH}, qui peut \^etre r\'ef\'erenc\'ee dans la requ\^ete principale.
Elle am\'eliore la lisibilit\'e et permet la r\'ecursion.
\end{definition}

\begin{codeblock}[title=\textbf{CTE simple}]
\begin{minted}{sql}
WITH Moyennes AS (
    SELECT id_etudiant, AVG(note) AS moy
    FROM Inscriptions
    GROUP BY id_etudiant
)
SELECT e.nom, m.moy
FROM Etudiants e
JOIN Moyennes m ON e.id_etudiant = m.id_etudiant
WHERE m.moy > 14
ORDER BY m.moy DESC;
\end{minted}
\end{codeblock}

\begin{codeblock}[title=\textbf{CTE r\'ecursive --- hi\'erarchie}]
\begin{minted}{sql}
-- Table de categories avec parent
CREATE TABLE Categories (
    id INTEGER PRIMARY KEY,
    nom TEXT,
    parent_id INTEGER REFERENCES Categories(id)
);

-- Parcours recursif de la hierarchie
WITH RECURSIVE Arbre AS (
    -- Cas de base : racines
    SELECT id, nom, parent_id, 0 AS niveau
    FROM Categories WHERE parent_id IS NULL

    UNION ALL

    -- Cas recursif
    SELECT c.id, c.nom, c.parent_id, a.niveau + 1
    FROM Categories c
    JOIN Arbre a ON c.parent_id = a.id
)
SELECT * FROM Arbre ORDER BY niveau, nom;
\end{minted}
\end{codeblock}

\section{Int\'egration Python}

\begin{codeblock}[title=\textbf{Jointures et sous-requ\^etes avec Python}]
\begin{minted}{python}
import sqlite3

conn = sqlite3.connect('universite.db')
cur = conn.cursor()

# Jointure : etudiants et leurs cours
cur.execute('''
    SELECT e.nom, c.intitule, i.note
    FROM Etudiants e
    JOIN Inscriptions i ON e.id_etudiant = i.id_etudiant
    JOIN Cours c ON i.id_cours = c.id_cours
    WHERE i.note > ?
    ORDER BY i.note DESC
''', (14,))

print("Etudiants avec note > 14 :")
for nom, cours, note in cur.fetchall():
    print(f"  {nom:10} | {cours:20} | {note:.1f}")

# Sous-requete : etudiants au-dessus de la moyenne
cur.execute('''
    SELECT nom, prenom FROM Etudiants
    WHERE id_etudiant IN (
        SELECT id_etudiant FROM Inscriptions
        WHERE note > (SELECT AVG(note) FROM Inscriptions)
    )
''')
print("\nAu-dessus de la moyenne :", cur.fetchall())

conn.close()
\end{minted}
\end{codeblock}

\section*{Exercices}

\begin{exercice}
\'Ecrivez une requ\^ete qui affiche le nom de chaque professeur et le nombre
de cours qu'il enseigne. Incluez les professeurs qui n'enseignent aucun cours.
\end{exercice}

\begin{exercice}
Trouvez les \'etudiants qui sont inscrits \`a tous les cours du professeur
dont l'\texttt{id\_prof} est 2. Utilisez une sous-requ\^ete avec \texttt{NOT EXISTS}.
\end{exercice}

\begin{exercice}
\'Ecrivez la m\^eme requ\^ete que l'exercice pr\'ec\'edent en utilisant
\texttt{GROUP BY} et \texttt{HAVING} avec \texttt{COUNT}.
Comparez les deux approches.
\end{exercice}

\begin{exercice}
Trouvez les paires d'\'etudiants qui ont au moins deux cours en commun.
Affichez les noms des deux \'etudiants et le nombre de cours partag\'es.
\end{exercice}

\begin{exercice}
En utilisant une CTE r\'ecursive, g\'en\'erez la table des nombres de 1 \`a 20
et affichez leur factorielle.
\end{exercice}
